%
% Introduction
%
\section{Introduction}
\label{sec:Introduction}

\subsection{Background and Motivation}
Quantum computing represents a paradigm shift in computational science, offering
exponential speedups for specific classes of problems, including quantum chemistry simulation \cite{Kandala2017},
cryptography \cite{cain2026shorsalgorithmpossible10000}, and
large-scale optimization \cite{Ye2025} in many industry and scientific areas. At the heart of this physical realization
are quantum bits (qubits). Among the various physical implementations, superconducting qubits---particularly the transmon
architecture---have established themselves as a leading foundational platform \cite{Acharya2025}.
Their macroscopic nature allows for strong coupling to microwave control fields and flexible
lithographic fabrication, making them highly scalable \cite{Devoret2013}. However, realizing fault-tolerant quantum computation requires
extremely high-fidelity quantum operations. Current Noisy Intermediate-Scale Quantum (NISQ) devices are fundamentally
limited by decoherence and operational errors. For superconducting qubits, single- and two-qubit gate fidelities are
highly sensitive to the precise shape of the microwave control pulses used to drive the transitions. Imperfect control
pulses lead to phase errors, over-rotation, and leakage out of the computational subspace, severely degrading the
overall performance of the quantum processor.

\subsection{Problem Statement}
The transmon qubit is inherently a weakly anharmonic oscillator. Because the energy difference between the
$\vert{}0\rangle \leftrightarrow \vert{}1\rangle$ transition and the $\vert{}1\rangle \leftrightarrow \vert{}2\rangle$
transition is relatively small, applying fast, high-amplitude microwave pulses to minimize decoherence time inevitably
results in spectral leakage into higher, non-computational energy states.
The Derivative Removal by Adiabatic Gate (DRAG) technique is a highly successful methodology for suppressing this
leakage. However, the demand for even higher fidelities necessitates further refinement of the underlying waveforms used
within the DRAG framework.
Because standard Gaussian or square base pulses can be spectrally broad or susceptible to noise, there is a critical
need to integrate mathematically optimal pulse shapes into the DRAG protocol. Specifically, by
actively suppressing the amplitude of these optimized waveforms, this combined approach can strictly confine the energy
spectrum while maintaining fast gate times. Implementing this requires robust calibration frameworks to validate these
amplitude-suppressed pulses against realistic Hamiltonian dynamics.
\subsection{Objectives and Contributions}
The primary objective of this thesis is to investigate, optimize, and evaluate advanced gate pulse waveforms for
superconducting transmon qubits, both one qubit and two qubit operations. By uncovering the underlying spectral
properties of established control techniques and introducing novel waveform designs, this work minimizes leakage errors
while accelerating gate execution times. The specific original contributions of this thesis are detailed as follows:

\noindent\textbf{Single-Qubit Gate Innovations}
\begin{description}
  \item[Spectral Formulation of Recursive DRAG ] Successfully demonstrated that the condition equations for DRAG-1 and DRAG-2
        pulses can be rigorously defined as a zero Fourier transform requirement. This connects the superconducting DRAG
        formalism to analogous mathematical conditions frequently observed in spin qubit research \cite{freeman,hoult1979,vold1968,warren1984}, providing a deeper physical
        intuition for recursive DRAG schemes.


  \item[Novel Pulse Design Methodology ] Leveraged this newly established Fourier transform formulation to develop a
        novel analytical pulse design approach. This method bypasses conventional constraints and provides an improved
        control over the resulting pulse shape, successfully achieving shorter
        single-qubit gate durations without compromising operational fidelity.
\end{description}

\noindent\textbf{Two-Qubit Gate Innovations}

\begin{description}
  \item[Anti-Symmetric Chebyshev Pulses ] Discovered and proposed a novel control waveform, designated as the anti-symmetric
        Chebyshev pulse, tailored to optimize two-qubit gate operations.

  \item[Targeted Spectral Suppression ] Adapted and extended the traditional Slepian pulse optimization framework by applying a
        technique specifically designed to minimize the Fourier transform amplitude within a precisely targeted frequency range,
        thereby optimizing the waveform for two-qubit leakage suppression.
\end{description}

\subsection{Thesis Organization}
The remainder of this thesis is structured to guide the
reader from the fundamental physics of the system through the mathematical modeling and finally to the numerical
results of the optimized waveforms.

